%!TEX root = ../../exa-ma-deliverables.tex
% Sources: external partners workshop and January 2026 General Assembly slides,
% Exa-MA partner register, and the June 2026 ANR review.
\begin{table}[!htb]
\centering
\caption{Doctoral projects co-funded by Exa-MA and industrial partners.}
\label{tab:exama-industrial-phds-2026}
\small
\setlength{\tabcolsep}{4pt}
\renewcommand{\arraystretch}{1.12}
\begin{tabular}{@{}
  >{\raggedright\arraybackslash}p{0.285\linewidth}
  >{\raggedright\arraybackslash}p{0.255\linewidth}
  >{\raggedright\arraybackslash}p{0.415\linewidth}@{}}
\toprule
\textbf{Partners and host team} & \textbf{Funding} & \textbf{Research}\\
\midrule
IFPEN and ONERA \newline Inria Alpines &
50\% Exa-MA, 25\% IFPEN and 25\% ONERA &
Learning-based solvers and preconditioners for sparse linear systems.\\[2pt]
ONERA and Safran \newline École polytechnique &
50\% Exa-MA, 50\% ONERA and Safran jointly &
High-order adaptive time coupling for multiphysics simulations.\\[2pt]
EDF \newline Sorbonne Université &
Exa-MA and EDF &
Scalable preconditioners for nonsymmetric saddle-point problems in WP3.\\[2pt]
TotalEnergies \newline Inria Makutu, Pau &
50\% Exa-MA, 50\% TotalEnergies &
Modular HPC software for solving PDEs, including AI-assisted design of numerical schemes.\\
\bottomrule
\end{tabular}
\end{table}
